\documentclass[11pt]{article}
\usepackage[a4paper,margin=1in]{geometry}
\usepackage{fontspec}
\usepackage[boldfont=false]{xeCJK}
\setCJKmainfont{FandolSong-Regular.otf}
\usepackage{amsmath}
\usepackage{amssymb}
\usepackage{array}
\usepackage{booktabs}
\usepackage{graphicx}
\usepackage{adjustbox}
\usepackage{enumitem}
\setlist[itemize]{topsep=2pt,partopsep=0pt,parsep=0pt,itemsep=2pt,leftmargin=1.4em}
\setlist[enumerate]{topsep=2pt,partopsep=0pt,parsep=0pt,itemsep=2pt,leftmargin=1.6em}
\usepackage{parskip}
\usepackage{fvextra}
\fvset{breaklines=true,breakanywhere=true,fontsize=\small}
\setlength{\parindent}{0pt}

\begin{document}

\begin{center}
  {\LARGE\bfseries Marketing}\\[0.35em]
  {\large Igcse Business Studies}
\end{center}
\vspace{0.6em}

\section*{What marketing does}

\textbf{Marketing} 营销 is the work of finding out what \textbf{customers} 顾客 want, then making, pricing, selling and promoting products to meet those wants at a profit.

Marketing has several roles:

\begin{itemize}
\item find out and satisfy what customers want,

\item keep and increase \textbf{sales} 销售额 and \textbf{market share} 市场份额,

\item build \textbf{customer loyalty} 顾客忠诚度 so buyers come back again.

\end{itemize}

\subsection*{Why customers matter}

Customers bring in the money, so keeping them is vital. Losing customers is costly: the business loses their future spending, and it must spend a lot on \textbf{advertising} 广告 to win new ones. It is usually far cheaper to keep a customer than to find a new one.

\subsection*{Niche and mass markets}

A \textbf{market} 市场 is all the buyers and sellers of a product.

\begin{itemize}
\item A \textbf{niche market} 利基市场 is a small part of a market with special needs, for example left-handed tools.

\item A \textbf{mass market} 大众市场 is a large market with many similar buyers, for example soft drinks.

\end{itemize}

\subsection*{Market segmentation}

\textbf{Market segmentation} 市场细分 means splitting a market into groups of buyers who share something, such as age, income, gender or location. Its benefits:

\begin{itemize}
\item the business can design a product to fit each group,

\item it can aim its advertising at the right people,

\item it can spot gaps in the market.

\end{itemize}

\section*{Market research}

\textbf{Market research} 市场调研 is collecting information about customers, rivals and the market, to help make good decisions.

\subsection*{Primary and secondary research}

\begin{itemize}
\item \textbf{Primary research} 一手调研 (also called field research) collects new information directly, for the first time. Methods include \textbf{questionnaires} 问卷, interviews, surveys and tests. It is up to date and fits the exact need, but costs time and money.

\item \textbf{Secondary research} 二手调研 (also called desk research) uses information that already exists, such as government reports, websites and past sales records. It is cheap and quick, but may be old or may not fit the need.

\end{itemize}

\subsection*{Sampling and reliable data}

A business cannot ask everyone, so it asks a small group, called a sample. \textbf{Sampling} 抽样 must be done with care. If the sample is too small or not typical, the \textbf{data} 数据 will not be reliable, and decisions based on it may be wrong.

\subsection*{Using research data}

Research results are shown in tables, charts and graphs. You should be able to read them — for example, find the most popular product, or see how sales change over time — and use them to support a decision.

\section*{The marketing mix}

The \textbf{marketing mix} 营销组合 is the four main choices a business makes about a product. They are often called the \textbf{four Ps}: product, price, place and promotion.

\subsection*{Product}

The \textbf{product} 产品 must meet customers' needs. Two ideas matter here.

The \textbf{product life cycle} 产品生命周期 describes the sales of a product over time, in four stages: introduction, growth, maturity and decline. When sales begin to fall, a business can use an \textbf{extension strategy} 延长策略 — such as a new design, a new advert, or selling in a new market — to keep the product selling.

A strong \textbf{brand} 品牌 (a name, logo and image that customers recognise) and good \textbf{packaging} 包装 (the wrapping that protects the product and attracts buyers) both help a product sell. Building a brand is called branding.

\subsection*{Price}

The \textbf{price} 价格 is what the customer pays. A business can choose from several pricing methods.

\begin{center}
\adjustbox{max width=\linewidth}{%
\begin{tabular}{ll}
\toprule
\textbf{Method} & \textbf{How it works} \\
\midrule
\textbf{cost-plus pricing} 成本加成定价 & add a fixed share of profit on top of the cost of making the item \\
\textbf{penetration pricing} 渗透定价 & set a low price to enter a market and win customers fast \\
\textbf{price skimming} 撇脂定价 & set a high price at first for a new, special product \\
\textbf{competitive pricing} 竞争定价 & set a price close to rivals' prices \\
\textbf{promotional pricing} 促销定价 & cut the price for a short time to lift sales \\
\textbf{dynamic pricing} 动态定价 & change the price automatically as demand changes, as airlines do \\
\bottomrule
\end{tabular}}
\end{center}

The right price also depends on the cost of making the product, on what customers will pay, on rivals' prices, and on the product's stage in its life cycle.

\subsection*{Place}

\textbf{Place} 地点 is about how and where the product is sold, so it reaches the customer at the right time. The route a product takes from maker to customer is called a \textbf{channel of distribution} 分销渠道.

Many channels use a \textbf{wholesaler} 批发商 (buys large amounts from makers and sells smaller amounts to shops) and a \textbf{retailer} 零售商 (a shop that sells to the final customer). A business can also sell \textbf{directly} 直销 to customers, or online through \textbf{e-commerce} 电子商务.

\subsection*{Promotion}

\textbf{Promotion} 推广 means telling customers about a product and persuading them to buy it. It includes:

\begin{itemize}
\item \textbf{advertising} — paid messages on TV, online, on posters and so on,

\item \textbf{sales promotion} 促销 — short-term offers like discounts, free gifts, or "buy one get one free".

\end{itemize}

Technology now plays a big part. The \textbf{internet} 互联网 and \textbf{social media} 社交媒体 let a business reach many people cheaply, target the right customers, and sell all over the world.

\section*{Marketing strategy}

A \textbf{marketing strategy} 营销策略 is a plan that uses the four Ps together to reach the business's \textbf{objectives} 目标. The four Ps must fit each other and fit the target customers. For example, a high-price luxury product needs high-quality packaging, expensive shops and image-based adverts.

\subsection*{Legal controls on marketing}

Many countries have \textbf{legal controls} 法律管制 that protect buyers:

\begin{itemize}
\item \textbf{consumer protection} 消费者保护 laws stop the sale of unsafe or faulty goods,

\item rules ban \textbf{misleading} 误导 promotion — adverts that lie about what a product can do.

\end{itemize}

\subsection*{Selling in other countries}

A business can grow by entering an \textbf{international market} 国际市场 in another country.

\begin{itemize}
\item Opportunities: many more customers, higher sales, and spreading risk across markets.

\item Problems: different languages and tastes, higher transport costs, foreign laws and \textbf{tariffs} 关税 (taxes on imports), and strong local \textbf{competition} 竞争.

\end{itemize}


\end{document}
